%% kapa: exam=1 round=33 year=2022 period=2 form=A subject=appraisal_law subject_ko=감정평가관계법규 seq=22 number=22
%% source: Q-Net 감정평가사 기출문제 (cst00309) artlSeq=5209997 file=raw/1008/5209997/2022년도 제33회 감정평가사 1차시험 2교시(감정평가 관계법규, 회계학).pdf
%% conversion: kordoc 4.24.51 -> tools/gen1.py (auto); review=proofread-sonnet; flags=none
\begin{question}{22}
국유재산법령상 일반재산의 처분가격에 관한 설명으로 옳은 것은?
\begin{choices}
  \item 증권을 처분할 때에는 시가를 고려하여 예정가격을 결정하여야 한다.
  \item 공공기관에 일반재산을 처분하는 경우에는 두 개의 감정평가법인등의 평가액을 산술평균한 금액을 예정가격으로 하여야 한다.
  \item 감정평가법인등의 평가액은 평가일부터 2년까지 적용할 수 있다.
  \item 국가가 보존·활용할 필요가 없고 대부·매각이나 교환이 곤란하여 일반재산을 양여하는 경우에는 대장가격을 재산가격으로 한다.
  \item 일단(一團)의 토지 대장가격이 3천만원 이하인 국유지를 경쟁입찰의 방법으로 처분하는 경우에는 해당 국유지의 개별공시지가를 예정가격으로 할 수 있다.
\end{choices}
\end{question}
\answer{4}
